% Part: second-order-logic
% Chapter: syntax-and-semantics
% Section: introduction

\documentclass[../../../include/open-logic-section]{subfiles}

\begin{document}

\olfileid{sol}{syn}{int}

\olsection{પ્રસ્તાવના}

પ્રથમ-ક્રમ તર્કશાસ્ત્રમાં આપેલી ભાષાનાં અતાર્કિક સંકેતો, એટલે કે તેનાં
!!{constant}s, !!{function}s અને !!{predicate}s, તાર્કિક સંકેતો સાથે
જોડીને આપણે પ્રથમ-ક્રમ !!{structure}s વિશેનાં વિધાનો વ્યક્ત કરીએ છીએ.
આ માટે સંતોષનો ખ્યાલ વપરાય છે. તે !!a{structure}~$\Struct{M}$,
ચલ-નિયુક્તિ~$s$ અને !!a{formula}~$!A$ને જોડે છે: $\Sat{M}{!A}[s]$
ત્યારે અને માત્ર ત્યારે જ સાચું છે જ્યારે~$!A$નાં !!{constant}s,
!!{function}s અને !!{predicate}sનું~$\Struct{M}$ મુજબ અર્થઘટન તથા
તેના મુક્ત ચલોનું~$s$ મુજબ અર્થઘટન કરીએ ત્યારે તે જે કહે છે તે
સાચું હોય. !!{identity}~$\eq$નું અર્થઘટન, તેમજ~$\lforall$ અને
$\lexists$નું અર્થઘટન, $\Sat{M}{!A}[s]$ની વ્યાખ્યામાં જ નિહિત છે.
સમાનતા હંમેશાં સંરચનાના !!{domain}~$\Domain{M}$ પરના અભિન્નતા
સંબંધ તરીકે અર્થઘટિત થાય છે અને પરિમાણકો હંમેશાં આખા !!{domain} પર
વ્યાપે છે. પરંતુ અગત્યની મર્યાદા એ છે કે પરિમાણીકરણ માત્ર
!!{domain}ના ઘટકો પર જ થઈ શકે છે; તેથી પરિમાણક પછી માત્ર વસ્તુ-!!{variable}s
આવી શકે છે.

દ્વિતીય-ક્રમ તર્કશાસ્ત્રમાં ભાષા અને સંતોષની વ્યાખ્યા બંનેને મુક્ત અને
બદ્ધ વિધેય-ચલો તથા સંબંધ-ચલો અને તેમના પરના પરિમાણીકરણનો સમાવેશ થાય
એ રીતે વિસ્તૃત કરીએ છીએ. વસ્તુ-ચલોનો !!{constant}s સાથે જેવો સંબંધ છે,
એવો જ આ ચલોનો !!{function}s અને !!{predicate}s સાથે છે. દ્વિતીય-ક્રમ
પદો અને !!{formula}sની રચનામાં તેઓ એ જ પ્રકારની ભૂમિકા ભજવે છે અને
તેમના પરનું પરિમાણીકરણ પણ સમાન રીતે સંભાળવામાં આવે છે. \emph{પ્રમાણભૂત}
અર્થવિજ્ઞાનમાં દ્વિતીય-ક્રમ પરિમાણકો યોગ્ય પ્રકારની બધી શક્ય વસ્તુઓ
પર વ્યાપે છે: વિધેય-ચલો માટે $\Domain{M}$માંથી~$\Domain{M}$માં જતાં
$n$-સ્થાની વિધેયો અને વિધેયધર્મ-ચલો માટે $n$-સ્થાની સંબંધો.
દાખલા તરીકે,
$\lforall[\Obj{v_0}][(\Obj{P^1_0}(\Obj{v_0}) \lor \lnot
  \Obj{P^1_0}(\Obj{v_0}))]$ પ્રથમ-ક્રમ અને દ્વિતીય-ક્રમ બંને
તર્કશાસ્ત્રમાં સૂત્ર છે. દ્વિતીય-ક્રમ તર્કશાસ્ત્રમાં વળી
$\lforall[\Obj{V^1_0}][\lforall[\Obj{v_0}][(\Obj{V^1_0}(\Obj{v_0}) \lor
    \lnot \Obj{V^1_0}(\Obj{v_0}))]]$ અને
$\lexists[\Obj{V^1_0}][\lforall[\Obj{v_0}][(\Obj{V^1_0}(\Obj{v_0}) \lor
    \lnot \Obj{V^1_0}(\Obj{v_0}))]]$ પણ વિચારી શકાય. તેમાં કોઈ મુક્ત
ચલ નથી, તેથી તે દ્વિતીય-ક્રમ તર્કશાસ્ત્રનાં !!{sentence}s છે. અહીં
$\Obj{V^1_0}$ દ્વિતીય-ક્રમનો $1$-સ્થાની વિધેયધર્મ-ચલ છે. $\Obj{V^1_0}$નું
સંભવિત અર્થઘટન~$\Obj{P^1_0}$ જેવા $1$-સ્થાની !!{predicate}ને આપી
શકાતા અર્થઘટન જેવું જ છે, એટલે~$\Domain{M}$ના ઉપગણો. તેથી આવા ચલો
પર પરિમાણીકરણનો અર્થ એ થાય છે કે
$\lforall[\Obj{v_0}][(\Obj{V^1_0}(\Obj{v_0}) \lor \lnot
  \Obj{V^1_0}(\Obj{v_0}))]$ એ~$\Obj{V^1_0}$ને~$\Domain{M}$નો કોઈપણ
ઉપગણ આપીએ ત્યારે સાચું હોય છે, અથવા ઓછામાં ઓછો એક એવો ઉપગણ આપીએ
ત્યારે સાચું હોય છે. દરેક ગણ આપેલા ઘટકને કાં તો સમાવે છે અથવા નથી
સમાવતો, તેથી બંને દ્વિતીય-ક્રમ વાક્યો કોઈપણ !!{structure}માં સાચાં છે.
\sourcecorrection{OLSOL-001}{સ્થિર અંગ્રેજી મૂળના છેલ્લા સૂત્રમાં
$\Obj v_0$ અને પછી $v_0$ લખાયું છે. ઉપરના બદ્ધ વસ્તુ-ચલ માટેના
એકસરખા $\Obj{v_0}$ સંકેતથી બંને સ્થાનો સ્પષ્ટ કર્યાં છે.}
\end{document}
